% =============================================================================
%  Fundamentos compartilhados --- Módulo 08: arquitetura de software
%  Incluído pelas trilhas Kotlin e Flutter com % =============================================================================
%  Fundamentos compartilhados --- Módulo 08: arquitetura de software
%  Incluído pelas trilhas Kotlin e Flutter com % =============================================================================
%  Fundamentos compartilhados --- Módulo 08: arquitetura de software
%  Incluído pelas trilhas Kotlin e Flutter com % =============================================================================
%  Fundamentos compartilhados --- Módulo 08: arquitetura de software
%  Incluído pelas trilhas Kotlin e Flutter com \input{../fundamentos/f08-arquitetura}
%  (coesão, acoplamento, SOLID e padrões de projeto estão no Módulo 03)
% =============================================================================
\section{Fundamentos: arquitetura de software}

\begin{frame}{O que é arquitetura de software}
  \begin{columns}[T]
    \begin{column}{0.52\textwidth}
      \begin{definicao}[Arquitetura]
        As grandes estruturas de um sistema: seus elementos, as relações entre eles e as decisões \textbf{difíceis de mudar depois}. Ela existe para proteger atributos de qualidade.
      \end{definicao}
      \scriptsize
      \begin{tabularx}{\linewidth}{@{}l L@{}}
        \toprule
        \cabecalho{Atributo} & \cabecalho{Pergunta no LabMonitor} \\
        \midrule
        Manutenibilidade & trocar a API exige mexer nas telas? \\
        Testabilidade & dá para testar a regra de reserva sem um aparelho? \\
        Desempenho & a lista de leituras rola sem engasgar? \\
        Funcionamento \emph{offline} & o app mostra dados sem internet? \\
        Segurança & o \emph{token} fica protegido? \\
        \bottomrule
      \end{tabularx}
    \end{column}
    \begin{column}{0.46\textwidth}
      \begin{analogia}
        A planta de uma casa decide por onde passam os canos antes da primeira parede. Pintar um quarto não é arquitetura; mudar a cozinha de lugar é.
      \end{analogia}
      \begin{conceito}[Boa arquitetura em um app]
        Permite mudar a interface sem tocar nos dados, trocar a API sem tocar na interface e testar regras sem abrir o app.
      \end{conceito}
      \begin{dica}
        Coesão, acoplamento, SOLID e padrões de projeto foram vistos no Módulo 03: aqui eles ganham escala.
      \end{dica}
    \end{column}
  \end{columns}
\end{frame}

\begin{frame}{Padrões de apresentação: de MVC a MVI}
  \centering
  \begin{tikzpicture}[every node/.style={font=\tiny}]
    \draw[seta] (0,0) -- (13.4,0);
    \foreach \x/\t/\d in {1.2/MVC/{1979, Smalltalk\\ (Trygve Reenskaug)},4.6/MVP/{anos 1990 (Taligent);\\ comum no Android até 2017},8.1/MVVM/{2005, Microsoft (WPF);\\ padrão no Android e no Flutter},11.6/MVI/{anos 2010, inspirado\\ no Redux e no Cycle.js}}{
      \fill[ambar] (\x,0) circle (2pt);
      \node[font=\scriptsize\bfseries,text=azulpetroleo,above=1mm] at (\x,0) {\t};
      \node[rotulo,below=1mm] at (\x,0) {\d};
    }
  \end{tikzpicture}
  \raggedright
  \scriptsize
  \begin{tabularx}{\textwidth}{@{}l L L L@{}}
    \toprule
    \cabecalho{Padrão} & \cabecalho{Quem atualiza a tela} & \cabecalho{Ponto forte} & \cabecalho{Ponto fraco} \\
    \midrule
    MVC & o \emph{controller} manipula a \emph{view} & simples & no Android, a \cod{Activity} virava um \emph{controller} gigante \\
    MVP & o \emph{presenter} chama métodos de uma interface da \emph{view} & testável sem a tela & muito código de interface; guarda referência à \emph{view} \\
    MVVM & a \emph{view} observa o estado exposto pelo \emph{view model} & casa com UI declarativa e sobrevive à rotação & vários fluxos de estado podem se desencontrar \\
    MVI & estado único e imutável; cada intenção gera um novo estado & previsível e fácil de depurar & cerimônia demais para telas simples \\
    \bottomrule
  \end{tabularx}
  \normalsize
  \begin{conceito}[Na prática atual]
    Compose e Flutter usam MVVM com um estado único por tela (\cod{UiState}): uma mistura de MVVM com a disciplina do MVI.
  \end{conceito}
\end{frame}

\begin{frame}{Arquitetura limpa e arquitetura hexagonal}
  \begin{columns}[c]
    \begin{column}{0.46\textwidth}
      \centering
      \begin{tikzpicture}[every node/.style={font=\tiny,align=center}]
        \fill[cinzaclaro] (0,0) circle (2.15);
        \fill[cinzafundo,draw=azulpetroleo!40] (0,0) circle (1.6);
        \fill[ambarclaro!60,draw=ambar] (0,0) circle (1.05);
        \fill[azulpetroleo] (0,0) circle (0.55);
        \node[text=white,font=\tiny\bfseries] at (0,0) {Entidades};
        \node at (0,0.8) {Casos de uso};
        \node at (0,1.33) {Adaptadores};
        \node at (0,1.88) {Frameworks e drivers};
        \node[rotulo] at (0,-1.33) {\emph{view models}, repositórios};
        \node[rotulo] at (0,-1.88) {UI, banco, HTTP, BLE};
        \draw[setaambar] (2.6,-1.2) -- (1.2,-0.55);
        \node[rotulo,text width=1.5cm] at (3.05,-1.5) {dependências apontam para dentro};
      \end{tikzpicture}
    \end{column}
    \begin{column}{0.52\textwidth}
      \begin{definicao}[Regra da dependência]
        O código das camadas internas não conhece as externas: as regras de negócio não importam nada de Android, Flutter, HTTP ou SQL.
      \end{definicao}
      \begin{historia}
        Alistair Cockburn propôs a arquitetura hexagonal, de portas e adaptadores, em 2005. Robert C. Martin reuniu ideias parecidas na Arquitetura Limpa, em 2012.
      \end{historia}
      \begin{dica}[Na medida certa]
        Em apps pequenos, camadas de UI e de dados bastam. Crie a camada de domínio quando regras se repetirem ou precisarem de testes isolados, como recomendam o Android e o Flutter.
      \end{dica}
    \end{column}
  \end{columns}
\end{frame}

\begin{frame}[fragile]{Inversão de dependência na prática}
  \begin{columns}[T]
    \begin{column}{0.5\textwidth}
      \centering
      \begin{tikzpicture}[every node/.style={font=\tiny}]
        \node[rotulo,anchor=west] at (-1.1,0.55) {\textbf{sem inversão}};
        \node[caixaescura,font=\tiny,minimum width=2cm] (r1) at (0,0) {regras de leitura};
        \node[caixa,font=\tiny,minimum width=1.8cm] (b1) at (3.4,0) {\cod{FonteBle}};
        \draw[seta] (r1) -- node[rotulo,above]{conhece} (b1);
        \node[rotulo,anchor=west] at (-1.1,-0.75) {\textbf{com inversão}};
        \node[caixaescura,font=\tiny,minimum width=2cm] (r2) at (0,-1.3) {regras de leitura};
        \node[caixaambar,font=\tiny,minimum width=2.4cm] (i2) at (3.4,-1.3) {\emph{interface}\\ \cod{FonteLeituras}};
        \node[caixa,font=\tiny,minimum width=1.35cm] (b2) at (1.95,-2.75) {\cod{FonteBle}};
        \node[caixa,font=\tiny,minimum width=1.35cm] (m2) at (3.4,-2.75) {\cod{FonteMqtt}};
        \node[caixa,font=\tiny,minimum width=1.35cm] (f2) at (4.85,-2.75) {\cod{FonteFake}};
        \draw[seta] (r2) -- node[rotulo,above]{usa} (i2);
        \foreach \n in {b2,m2,f2}{\draw[setaambar] (\n.north) -- (\n.north |- i2.south);}
        \node[rotulo,below=1mm of m2] {implementam a interface};
      \end{tikzpicture}
\begin{codigo}[title=Dart]{Dart}
abstract interface class FonteLeituras {
  Stream<Leitura> leituras();
}
\end{codigo}
    \end{column}
    \begin{column}{0.48\textwidth}
      \begin{definicao}[Inversão de dependência]
        Módulos de alto nível não dependem de detalhes; ambos dependem de abstrações. Quem precisa define a interface; a camada de dados a implementa.
      \end{definicao}
\begin{codigo}[title=Kotlin]{Kotlin}
interface FonteLeituras {
  fun leituras(): Flow<Leitura>
}
\end{codigo}
      \begin{conceito}[O que se ganha]
        Trocar BLE por MQTT sem mexer nas regras, testar com uma fonte falsa e fazer a tela antes de o hardware existir.
      \end{conceito}
    \end{column}
  \end{columns}
\end{frame}

\begin{frame}{Injeção de dependência e localizador de serviços}
  \begin{columns}[T]
    \begin{column}{0.5\textwidth}
      \begin{definicao}[Inversão de controle]
        A classe não cria o que usa: alguém de fora monta os objetos e os entrega. É o princípio de Hollywood: ``não nos chame, nós chamamos você''.
      \end{definicao}
      \centering
      \begin{tikzpicture}[every node/.style={font=\tiny}]
        \node[caixaambarforte,font=\tiny\bfseries,minimum width=5.6cm] (raiz) at (0,0) {raiz de composição (\cod{Application}, \cod{main})};
        \node[caixa,font=\tiny,minimum width=1.7cm] (h) at (-1.9,-1) {cliente HTTP};
        \node[caixa,font=\tiny,minimum width=1.7cm] (b) at (0,-1) {banco};
        \node[caixapetroleo,font=\tiny,minimum width=1.9cm] (r) at (-0.95,-2) {repositório};
        \node[caixaescura,font=\tiny,minimum width=1.7cm] (v) at (1.9,-2) {\emph{view model}};
        \draw[setatrac] (raiz.south -| h) -- (h); \draw[setatrac] (raiz.south -| b) -- (b);
        \draw[seta] (h) -- (r); \draw[seta] (b) -- (r);
        \draw[seta] (r) -- node[rotulo,above]{injeta} (v);
        \node[rotulo] at (1.9,-1.05) {tracejado: cria};
      \end{tikzpicture}
      \raggedright
      \begin{historia}
        Martin Fowler batizou o padrão de injeção de dependência em 2004.
      \end{historia}
    \end{column}
    \begin{column}{0.48\textwidth}
      \scriptsize
      \begin{tabularx}{\linewidth}{@{}l L@{}}
        \toprule
        \cabecalho{Abordagem} & \cabecalho{Como funciona} \\
        \midrule
        Construtor & recebe tudo pronto; dependências explícitas e fáceis de trocar em testes \\
        Contêiner & uma biblioteca monta o grafo: Hilt, Koin; Provider, Riverpod \\
        Localizador & a classe pede a um registro global (\cod{get\_it}): prático, mas esconde dependências e só falha em execução \\
        \bottomrule
      \end{tabularx}
      \normalsize
      \begin{pordentro}[Tempo de vida (escopo)]
        \textbf{Único no app}: cliente HTTP, banco. \textbf{Por tela}: \emph{view model}. \textbf{Por uso}: objeto novo a cada pedido. Escopo errado causa vazamento ou estado compartilhado por engano.
      \end{pordentro}
    \end{column}
  \end{columns}
\end{frame}

\begin{frame}{Decisões de arquitetura: trocas e registros}
  \begin{columns}[T]
    \begin{column}{0.5\textwidth}
      \begin{conceito}[Toda decisão é uma troca]
        Mais camadas trazem isolamento e testabilidade, mas também mais arquivos e indireção. Pese o tamanho do time, a vida esperada do app e o custo de errar.
      \end{conceito}
      \scriptsize
      \begin{tabularx}{\linewidth}{@{}l L@{}}
        \toprule
        \cabecalho{Princípio} & \cabecalho{Significado} \\
        \midrule
        YAGNI & não construa o que ainda não é necessário \\
        KISS & prefira a solução mais simples que funciona \\
        DRY & cada conhecimento em um só lugar \\
        Regra de três & abstraia na terceira repetição, não na primeira \\
        \bottomrule
      \end{tabularx}
    \end{column}
    \begin{column}{0.48\textwidth}
      \begin{definicao}[ADR (\emph{Architecture Decision Record})]
        Texto curto, versionado com o código, que registra uma decisão: contexto, opções, decisão e consequências. Responde no futuro: ``por que fizemos assim?''.
      \end{definicao}
      \begin{exemplo}[ADR 007 --- Estado das telas]
        \textbf{Contexto}: leituras do sensor em tempo real. \textbf{Decisão}: um \cod{UiState} imutável por tela, exposto como fluxo observável. \textbf{Consequência}: telas testáveis sem interface.
      \end{exemplo}
      \begin{historia}
        Michael Nygard propôs os ADRs em 2011.
      \end{historia}
    \end{column}
  \end{columns}
\end{frame}

%  (coesão, acoplamento, SOLID e padrões de projeto estão no Módulo 03)
% =============================================================================
\section{Fundamentos: arquitetura de software}

\begin{frame}{O que é arquitetura de software}
  \begin{columns}[T]
    \begin{column}{0.52\textwidth}
      \begin{definicao}[Arquitetura]
        As grandes estruturas de um sistema: seus elementos, as relações entre eles e as decisões \textbf{difíceis de mudar depois}. Ela existe para proteger atributos de qualidade.
      \end{definicao}
      \scriptsize
      \begin{tabularx}{\linewidth}{@{}l L@{}}
        \toprule
        \cabecalho{Atributo} & \cabecalho{Pergunta no LabMonitor} \\
        \midrule
        Manutenibilidade & trocar a API exige mexer nas telas? \\
        Testabilidade & dá para testar a regra de reserva sem um aparelho? \\
        Desempenho & a lista de leituras rola sem engasgar? \\
        Funcionamento \emph{offline} & o app mostra dados sem internet? \\
        Segurança & o \emph{token} fica protegido? \\
        \bottomrule
      \end{tabularx}
    \end{column}
    \begin{column}{0.46\textwidth}
      \begin{analogia}
        A planta de uma casa decide por onde passam os canos antes da primeira parede. Pintar um quarto não é arquitetura; mudar a cozinha de lugar é.
      \end{analogia}
      \begin{conceito}[Boa arquitetura em um app]
        Permite mudar a interface sem tocar nos dados, trocar a API sem tocar na interface e testar regras sem abrir o app.
      \end{conceito}
      \begin{dica}
        Coesão, acoplamento, SOLID e padrões de projeto foram vistos no Módulo 03: aqui eles ganham escala.
      \end{dica}
    \end{column}
  \end{columns}
\end{frame}

\begin{frame}{Padrões de apresentação: de MVC a MVI}
  \centering
  \begin{tikzpicture}[every node/.style={font=\tiny}]
    \draw[seta] (0,0) -- (13.4,0);
    \foreach \x/\t/\d in {1.2/MVC/{1979, Smalltalk\\ (Trygve Reenskaug)},4.6/MVP/{anos 1990 (Taligent);\\ comum no Android até 2017},8.1/MVVM/{2005, Microsoft (WPF);\\ padrão no Android e no Flutter},11.6/MVI/{anos 2010, inspirado\\ no Redux e no Cycle.js}}{
      \fill[ambar] (\x,0) circle (2pt);
      \node[font=\scriptsize\bfseries,text=azulpetroleo,above=1mm] at (\x,0) {\t};
      \node[rotulo,below=1mm] at (\x,0) {\d};
    }
  \end{tikzpicture}
  \raggedright
  \scriptsize
  \begin{tabularx}{\textwidth}{@{}l L L L@{}}
    \toprule
    \cabecalho{Padrão} & \cabecalho{Quem atualiza a tela} & \cabecalho{Ponto forte} & \cabecalho{Ponto fraco} \\
    \midrule
    MVC & o \emph{controller} manipula a \emph{view} & simples & no Android, a \cod{Activity} virava um \emph{controller} gigante \\
    MVP & o \emph{presenter} chama métodos de uma interface da \emph{view} & testável sem a tela & muito código de interface; guarda referência à \emph{view} \\
    MVVM & a \emph{view} observa o estado exposto pelo \emph{view model} & casa com UI declarativa e sobrevive à rotação & vários fluxos de estado podem se desencontrar \\
    MVI & estado único e imutável; cada intenção gera um novo estado & previsível e fácil de depurar & cerimônia demais para telas simples \\
    \bottomrule
  \end{tabularx}
  \normalsize
  \begin{conceito}[Na prática atual]
    Compose e Flutter usam MVVM com um estado único por tela (\cod{UiState}): uma mistura de MVVM com a disciplina do MVI.
  \end{conceito}
\end{frame}

\begin{frame}{Arquitetura limpa e arquitetura hexagonal}
  \begin{columns}[c]
    \begin{column}{0.46\textwidth}
      \centering
      \begin{tikzpicture}[every node/.style={font=\tiny,align=center}]
        \fill[cinzaclaro] (0,0) circle (2.15);
        \fill[cinzafundo,draw=azulpetroleo!40] (0,0) circle (1.6);
        \fill[ambarclaro!60,draw=ambar] (0,0) circle (1.05);
        \fill[azulpetroleo] (0,0) circle (0.55);
        \node[text=white,font=\tiny\bfseries] at (0,0) {Entidades};
        \node at (0,0.8) {Casos de uso};
        \node at (0,1.33) {Adaptadores};
        \node at (0,1.88) {Frameworks e drivers};
        \node[rotulo] at (0,-1.33) {\emph{view models}, repositórios};
        \node[rotulo] at (0,-1.88) {UI, banco, HTTP, BLE};
        \draw[setaambar] (2.6,-1.2) -- (1.2,-0.55);
        \node[rotulo,text width=1.5cm] at (3.05,-1.5) {dependências apontam para dentro};
      \end{tikzpicture}
    \end{column}
    \begin{column}{0.52\textwidth}
      \begin{definicao}[Regra da dependência]
        O código das camadas internas não conhece as externas: as regras de negócio não importam nada de Android, Flutter, HTTP ou SQL.
      \end{definicao}
      \begin{historia}
        Alistair Cockburn propôs a arquitetura hexagonal, de portas e adaptadores, em 2005. Robert C. Martin reuniu ideias parecidas na Arquitetura Limpa, em 2012.
      \end{historia}
      \begin{dica}[Na medida certa]
        Em apps pequenos, camadas de UI e de dados bastam. Crie a camada de domínio quando regras se repetirem ou precisarem de testes isolados, como recomendam o Android e o Flutter.
      \end{dica}
    \end{column}
  \end{columns}
\end{frame}

\begin{frame}[fragile]{Inversão de dependência na prática}
  \begin{columns}[T]
    \begin{column}{0.5\textwidth}
      \centering
      \begin{tikzpicture}[every node/.style={font=\tiny}]
        \node[rotulo,anchor=west] at (-1.1,0.55) {\textbf{sem inversão}};
        \node[caixaescura,font=\tiny,minimum width=2cm] (r1) at (0,0) {regras de leitura};
        \node[caixa,font=\tiny,minimum width=1.8cm] (b1) at (3.4,0) {\cod{FonteBle}};
        \draw[seta] (r1) -- node[rotulo,above]{conhece} (b1);
        \node[rotulo,anchor=west] at (-1.1,-0.75) {\textbf{com inversão}};
        \node[caixaescura,font=\tiny,minimum width=2cm] (r2) at (0,-1.3) {regras de leitura};
        \node[caixaambar,font=\tiny,minimum width=2.4cm] (i2) at (3.4,-1.3) {\emph{interface}\\ \cod{FonteLeituras}};
        \node[caixa,font=\tiny,minimum width=1.35cm] (b2) at (1.95,-2.75) {\cod{FonteBle}};
        \node[caixa,font=\tiny,minimum width=1.35cm] (m2) at (3.4,-2.75) {\cod{FonteMqtt}};
        \node[caixa,font=\tiny,minimum width=1.35cm] (f2) at (4.85,-2.75) {\cod{FonteFake}};
        \draw[seta] (r2) -- node[rotulo,above]{usa} (i2);
        \foreach \n in {b2,m2,f2}{\draw[setaambar] (\n.north) -- (\n.north |- i2.south);}
        \node[rotulo,below=1mm of m2] {implementam a interface};
      \end{tikzpicture}
\begin{codigo}[title=Dart]{Dart}
abstract interface class FonteLeituras {
  Stream<Leitura> leituras();
}
\end{codigo}
    \end{column}
    \begin{column}{0.48\textwidth}
      \begin{definicao}[Inversão de dependência]
        Módulos de alto nível não dependem de detalhes; ambos dependem de abstrações. Quem precisa define a interface; a camada de dados a implementa.
      \end{definicao}
\begin{codigo}[title=Kotlin]{Kotlin}
interface FonteLeituras {
  fun leituras(): Flow<Leitura>
}
\end{codigo}
      \begin{conceito}[O que se ganha]
        Trocar BLE por MQTT sem mexer nas regras, testar com uma fonte falsa e fazer a tela antes de o hardware existir.
      \end{conceito}
    \end{column}
  \end{columns}
\end{frame}

\begin{frame}{Injeção de dependência e localizador de serviços}
  \begin{columns}[T]
    \begin{column}{0.5\textwidth}
      \begin{definicao}[Inversão de controle]
        A classe não cria o que usa: alguém de fora monta os objetos e os entrega. É o princípio de Hollywood: ``não nos chame, nós chamamos você''.
      \end{definicao}
      \centering
      \begin{tikzpicture}[every node/.style={font=\tiny}]
        \node[caixaambarforte,font=\tiny\bfseries,minimum width=5.6cm] (raiz) at (0,0) {raiz de composição (\cod{Application}, \cod{main})};
        \node[caixa,font=\tiny,minimum width=1.7cm] (h) at (-1.9,-1) {cliente HTTP};
        \node[caixa,font=\tiny,minimum width=1.7cm] (b) at (0,-1) {banco};
        \node[caixapetroleo,font=\tiny,minimum width=1.9cm] (r) at (-0.95,-2) {repositório};
        \node[caixaescura,font=\tiny,minimum width=1.7cm] (v) at (1.9,-2) {\emph{view model}};
        \draw[setatrac] (raiz.south -| h) -- (h); \draw[setatrac] (raiz.south -| b) -- (b);
        \draw[seta] (h) -- (r); \draw[seta] (b) -- (r);
        \draw[seta] (r) -- node[rotulo,above]{injeta} (v);
        \node[rotulo] at (1.9,-1.05) {tracejado: cria};
      \end{tikzpicture}
      \raggedright
      \begin{historia}
        Martin Fowler batizou o padrão de injeção de dependência em 2004.
      \end{historia}
    \end{column}
    \begin{column}{0.48\textwidth}
      \scriptsize
      \begin{tabularx}{\linewidth}{@{}l L@{}}
        \toprule
        \cabecalho{Abordagem} & \cabecalho{Como funciona} \\
        \midrule
        Construtor & recebe tudo pronto; dependências explícitas e fáceis de trocar em testes \\
        Contêiner & uma biblioteca monta o grafo: Hilt, Koin; Provider, Riverpod \\
        Localizador & a classe pede a um registro global (\cod{get\_it}): prático, mas esconde dependências e só falha em execução \\
        \bottomrule
      \end{tabularx}
      \normalsize
      \begin{pordentro}[Tempo de vida (escopo)]
        \textbf{Único no app}: cliente HTTP, banco. \textbf{Por tela}: \emph{view model}. \textbf{Por uso}: objeto novo a cada pedido. Escopo errado causa vazamento ou estado compartilhado por engano.
      \end{pordentro}
    \end{column}
  \end{columns}
\end{frame}

\begin{frame}{Decisões de arquitetura: trocas e registros}
  \begin{columns}[T]
    \begin{column}{0.5\textwidth}
      \begin{conceito}[Toda decisão é uma troca]
        Mais camadas trazem isolamento e testabilidade, mas também mais arquivos e indireção. Pese o tamanho do time, a vida esperada do app e o custo de errar.
      \end{conceito}
      \scriptsize
      \begin{tabularx}{\linewidth}{@{}l L@{}}
        \toprule
        \cabecalho{Princípio} & \cabecalho{Significado} \\
        \midrule
        YAGNI & não construa o que ainda não é necessário \\
        KISS & prefira a solução mais simples que funciona \\
        DRY & cada conhecimento em um só lugar \\
        Regra de três & abstraia na terceira repetição, não na primeira \\
        \bottomrule
      \end{tabularx}
    \end{column}
    \begin{column}{0.48\textwidth}
      \begin{definicao}[ADR (\emph{Architecture Decision Record})]
        Texto curto, versionado com o código, que registra uma decisão: contexto, opções, decisão e consequências. Responde no futuro: ``por que fizemos assim?''.
      \end{definicao}
      \begin{exemplo}[ADR 007 --- Estado das telas]
        \textbf{Contexto}: leituras do sensor em tempo real. \textbf{Decisão}: um \cod{UiState} imutável por tela, exposto como fluxo observável. \textbf{Consequência}: telas testáveis sem interface.
      \end{exemplo}
      \begin{historia}
        Michael Nygard propôs os ADRs em 2011.
      \end{historia}
    \end{column}
  \end{columns}
\end{frame}

%  (coesão, acoplamento, SOLID e padrões de projeto estão no Módulo 03)
% =============================================================================
\section{Fundamentos: arquitetura de software}

\begin{frame}{O que é arquitetura de software}
  \begin{columns}[T]
    \begin{column}{0.52\textwidth}
      \begin{definicao}[Arquitetura]
        As grandes estruturas de um sistema: seus elementos, as relações entre eles e as decisões \textbf{difíceis de mudar depois}. Ela existe para proteger atributos de qualidade.
      \end{definicao}
      \scriptsize
      \begin{tabularx}{\linewidth}{@{}l L@{}}
        \toprule
        \cabecalho{Atributo} & \cabecalho{Pergunta no LabMonitor} \\
        \midrule
        Manutenibilidade & trocar a API exige mexer nas telas? \\
        Testabilidade & dá para testar a regra de reserva sem um aparelho? \\
        Desempenho & a lista de leituras rola sem engasgar? \\
        Funcionamento \emph{offline} & o app mostra dados sem internet? \\
        Segurança & o \emph{token} fica protegido? \\
        \bottomrule
      \end{tabularx}
    \end{column}
    \begin{column}{0.46\textwidth}
      \begin{analogia}
        A planta de uma casa decide por onde passam os canos antes da primeira parede. Pintar um quarto não é arquitetura; mudar a cozinha de lugar é.
      \end{analogia}
      \begin{conceito}[Boa arquitetura em um app]
        Permite mudar a interface sem tocar nos dados, trocar a API sem tocar na interface e testar regras sem abrir o app.
      \end{conceito}
      \begin{dica}
        Coesão, acoplamento, SOLID e padrões de projeto foram vistos no Módulo 03: aqui eles ganham escala.
      \end{dica}
    \end{column}
  \end{columns}
\end{frame}

\begin{frame}{Padrões de apresentação: de MVC a MVI}
  \centering
  \begin{tikzpicture}[every node/.style={font=\tiny}]
    \draw[seta] (0,0) -- (13.4,0);
    \foreach \x/\t/\d in {1.2/MVC/{1979, Smalltalk\\ (Trygve Reenskaug)},4.6/MVP/{anos 1990 (Taligent);\\ comum no Android até 2017},8.1/MVVM/{2005, Microsoft (WPF);\\ padrão no Android e no Flutter},11.6/MVI/{anos 2010, inspirado\\ no Redux e no Cycle.js}}{
      \fill[ambar] (\x,0) circle (2pt);
      \node[font=\scriptsize\bfseries,text=azulpetroleo,above=1mm] at (\x,0) {\t};
      \node[rotulo,below=1mm] at (\x,0) {\d};
    }
  \end{tikzpicture}
  \raggedright
  \scriptsize
  \begin{tabularx}{\textwidth}{@{}l L L L@{}}
    \toprule
    \cabecalho{Padrão} & \cabecalho{Quem atualiza a tela} & \cabecalho{Ponto forte} & \cabecalho{Ponto fraco} \\
    \midrule
    MVC & o \emph{controller} manipula a \emph{view} & simples & no Android, a \cod{Activity} virava um \emph{controller} gigante \\
    MVP & o \emph{presenter} chama métodos de uma interface da \emph{view} & testável sem a tela & muito código de interface; guarda referência à \emph{view} \\
    MVVM & a \emph{view} observa o estado exposto pelo \emph{view model} & casa com UI declarativa e sobrevive à rotação & vários fluxos de estado podem se desencontrar \\
    MVI & estado único e imutável; cada intenção gera um novo estado & previsível e fácil de depurar & cerimônia demais para telas simples \\
    \bottomrule
  \end{tabularx}
  \normalsize
  \begin{conceito}[Na prática atual]
    Compose e Flutter usam MVVM com um estado único por tela (\cod{UiState}): uma mistura de MVVM com a disciplina do MVI.
  \end{conceito}
\end{frame}

\begin{frame}{Arquitetura limpa e arquitetura hexagonal}
  \begin{columns}[c]
    \begin{column}{0.46\textwidth}
      \centering
      \begin{tikzpicture}[every node/.style={font=\tiny,align=center}]
        \fill[cinzaclaro] (0,0) circle (2.15);
        \fill[cinzafundo,draw=azulpetroleo!40] (0,0) circle (1.6);
        \fill[ambarclaro!60,draw=ambar] (0,0) circle (1.05);
        \fill[azulpetroleo] (0,0) circle (0.55);
        \node[text=white,font=\tiny\bfseries] at (0,0) {Entidades};
        \node at (0,0.8) {Casos de uso};
        \node at (0,1.33) {Adaptadores};
        \node at (0,1.88) {Frameworks e drivers};
        \node[rotulo] at (0,-1.33) {\emph{view models}, repositórios};
        \node[rotulo] at (0,-1.88) {UI, banco, HTTP, BLE};
        \draw[setaambar] (2.6,-1.2) -- (1.2,-0.55);
        \node[rotulo,text width=1.5cm] at (3.05,-1.5) {dependências apontam para dentro};
      \end{tikzpicture}
    \end{column}
    \begin{column}{0.52\textwidth}
      \begin{definicao}[Regra da dependência]
        O código das camadas internas não conhece as externas: as regras de negócio não importam nada de Android, Flutter, HTTP ou SQL.
      \end{definicao}
      \begin{historia}
        Alistair Cockburn propôs a arquitetura hexagonal, de portas e adaptadores, em 2005. Robert C. Martin reuniu ideias parecidas na Arquitetura Limpa, em 2012.
      \end{historia}
      \begin{dica}[Na medida certa]
        Em apps pequenos, camadas de UI e de dados bastam. Crie a camada de domínio quando regras se repetirem ou precisarem de testes isolados, como recomendam o Android e o Flutter.
      \end{dica}
    \end{column}
  \end{columns}
\end{frame}

\begin{frame}[fragile]{Inversão de dependência na prática}
  \begin{columns}[T]
    \begin{column}{0.5\textwidth}
      \centering
      \begin{tikzpicture}[every node/.style={font=\tiny}]
        \node[rotulo,anchor=west] at (-1.1,0.55) {\textbf{sem inversão}};
        \node[caixaescura,font=\tiny,minimum width=2cm] (r1) at (0,0) {regras de leitura};
        \node[caixa,font=\tiny,minimum width=1.8cm] (b1) at (3.4,0) {\cod{FonteBle}};
        \draw[seta] (r1) -- node[rotulo,above]{conhece} (b1);
        \node[rotulo,anchor=west] at (-1.1,-0.75) {\textbf{com inversão}};
        \node[caixaescura,font=\tiny,minimum width=2cm] (r2) at (0,-1.3) {regras de leitura};
        \node[caixaambar,font=\tiny,minimum width=2.4cm] (i2) at (3.4,-1.3) {\emph{interface}\\ \cod{FonteLeituras}};
        \node[caixa,font=\tiny,minimum width=1.35cm] (b2) at (1.95,-2.75) {\cod{FonteBle}};
        \node[caixa,font=\tiny,minimum width=1.35cm] (m2) at (3.4,-2.75) {\cod{FonteMqtt}};
        \node[caixa,font=\tiny,minimum width=1.35cm] (f2) at (4.85,-2.75) {\cod{FonteFake}};
        \draw[seta] (r2) -- node[rotulo,above]{usa} (i2);
        \foreach \n in {b2,m2,f2}{\draw[setaambar] (\n.north) -- (\n.north |- i2.south);}
        \node[rotulo,below=1mm of m2] {implementam a interface};
      \end{tikzpicture}
\begin{codigo}[title=Dart]{Dart}
abstract interface class FonteLeituras {
  Stream<Leitura> leituras();
}
\end{codigo}
    \end{column}
    \begin{column}{0.48\textwidth}
      \begin{definicao}[Inversão de dependência]
        Módulos de alto nível não dependem de detalhes; ambos dependem de abstrações. Quem precisa define a interface; a camada de dados a implementa.
      \end{definicao}
\begin{codigo}[title=Kotlin]{Kotlin}
interface FonteLeituras {
  fun leituras(): Flow<Leitura>
}
\end{codigo}
      \begin{conceito}[O que se ganha]
        Trocar BLE por MQTT sem mexer nas regras, testar com uma fonte falsa e fazer a tela antes de o hardware existir.
      \end{conceito}
    \end{column}
  \end{columns}
\end{frame}

\begin{frame}{Injeção de dependência e localizador de serviços}
  \begin{columns}[T]
    \begin{column}{0.5\textwidth}
      \begin{definicao}[Inversão de controle]
        A classe não cria o que usa: alguém de fora monta os objetos e os entrega. É o princípio de Hollywood: ``não nos chame, nós chamamos você''.
      \end{definicao}
      \centering
      \begin{tikzpicture}[every node/.style={font=\tiny}]
        \node[caixaambarforte,font=\tiny\bfseries,minimum width=5.6cm] (raiz) at (0,0) {raiz de composição (\cod{Application}, \cod{main})};
        \node[caixa,font=\tiny,minimum width=1.7cm] (h) at (-1.9,-1) {cliente HTTP};
        \node[caixa,font=\tiny,minimum width=1.7cm] (b) at (0,-1) {banco};
        \node[caixapetroleo,font=\tiny,minimum width=1.9cm] (r) at (-0.95,-2) {repositório};
        \node[caixaescura,font=\tiny,minimum width=1.7cm] (v) at (1.9,-2) {\emph{view model}};
        \draw[setatrac] (raiz.south -| h) -- (h); \draw[setatrac] (raiz.south -| b) -- (b);
        \draw[seta] (h) -- (r); \draw[seta] (b) -- (r);
        \draw[seta] (r) -- node[rotulo,above]{injeta} (v);
        \node[rotulo] at (1.9,-1.05) {tracejado: cria};
      \end{tikzpicture}
      \raggedright
      \begin{historia}
        Martin Fowler batizou o padrão de injeção de dependência em 2004.
      \end{historia}
    \end{column}
    \begin{column}{0.48\textwidth}
      \scriptsize
      \begin{tabularx}{\linewidth}{@{}l L@{}}
        \toprule
        \cabecalho{Abordagem} & \cabecalho{Como funciona} \\
        \midrule
        Construtor & recebe tudo pronto; dependências explícitas e fáceis de trocar em testes \\
        Contêiner & uma biblioteca monta o grafo: Hilt, Koin; Provider, Riverpod \\
        Localizador & a classe pede a um registro global (\cod{get\_it}): prático, mas esconde dependências e só falha em execução \\
        \bottomrule
      \end{tabularx}
      \normalsize
      \begin{pordentro}[Tempo de vida (escopo)]
        \textbf{Único no app}: cliente HTTP, banco. \textbf{Por tela}: \emph{view model}. \textbf{Por uso}: objeto novo a cada pedido. Escopo errado causa vazamento ou estado compartilhado por engano.
      \end{pordentro}
    \end{column}
  \end{columns}
\end{frame}

\begin{frame}{Decisões de arquitetura: trocas e registros}
  \begin{columns}[T]
    \begin{column}{0.5\textwidth}
      \begin{conceito}[Toda decisão é uma troca]
        Mais camadas trazem isolamento e testabilidade, mas também mais arquivos e indireção. Pese o tamanho do time, a vida esperada do app e o custo de errar.
      \end{conceito}
      \scriptsize
      \begin{tabularx}{\linewidth}{@{}l L@{}}
        \toprule
        \cabecalho{Princípio} & \cabecalho{Significado} \\
        \midrule
        YAGNI & não construa o que ainda não é necessário \\
        KISS & prefira a solução mais simples que funciona \\
        DRY & cada conhecimento em um só lugar \\
        Regra de três & abstraia na terceira repetição, não na primeira \\
        \bottomrule
      \end{tabularx}
    \end{column}
    \begin{column}{0.48\textwidth}
      \begin{definicao}[ADR (\emph{Architecture Decision Record})]
        Texto curto, versionado com o código, que registra uma decisão: contexto, opções, decisão e consequências. Responde no futuro: ``por que fizemos assim?''.
      \end{definicao}
      \begin{exemplo}[ADR 007 --- Estado das telas]
        \textbf{Contexto}: leituras do sensor em tempo real. \textbf{Decisão}: um \cod{UiState} imutável por tela, exposto como fluxo observável. \textbf{Consequência}: telas testáveis sem interface.
      \end{exemplo}
      \begin{historia}
        Michael Nygard propôs os ADRs em 2011.
      \end{historia}
    \end{column}
  \end{columns}
\end{frame}

%  (coesão, acoplamento, SOLID e padrões de projeto estão no Módulo 03)
% =============================================================================
\section{Fundamentos: arquitetura de software}

\begin{frame}{O que é arquitetura de software}
  \begin{columns}[T]
    \begin{column}{0.52\textwidth}
      \begin{definicao}[Arquitetura]
        As grandes estruturas de um sistema: seus elementos, as relações entre eles e as decisões \textbf{difíceis de mudar depois}. Ela existe para proteger atributos de qualidade.
      \end{definicao}
      \scriptsize
      \begin{tabularx}{\linewidth}{@{}l L@{}}
        \toprule
        \cabecalho{Atributo} & \cabecalho{Pergunta no LabMonitor} \\
        \midrule
        Manutenibilidade & trocar a API exige mexer nas telas? \\
        Testabilidade & dá para testar a regra de reserva sem um aparelho? \\
        Desempenho & a lista de leituras rola sem engasgar? \\
        Funcionamento \emph{offline} & o app mostra dados sem internet? \\
        Segurança & o \emph{token} fica protegido? \\
        \bottomrule
      \end{tabularx}
    \end{column}
    \begin{column}{0.46\textwidth}
      \begin{analogia}
        A planta de uma casa decide por onde passam os canos antes da primeira parede. Pintar um quarto não é arquitetura; mudar a cozinha de lugar é.
      \end{analogia}
      \begin{conceito}[Boa arquitetura em um app]
        Permite mudar a interface sem tocar nos dados, trocar a API sem tocar na interface e testar regras sem abrir o app.
      \end{conceito}
      \begin{dica}
        Coesão, acoplamento, SOLID e padrões de projeto foram vistos no Módulo 03: aqui eles ganham escala.
      \end{dica}
    \end{column}
  \end{columns}
\end{frame}

\begin{frame}{Padrões de apresentação: de MVC a MVI}
  \centering
  \begin{tikzpicture}[every node/.style={font=\tiny}]
    \draw[seta] (0,0) -- (13.4,0);
    \foreach \x/\t/\d in {1.2/MVC/{1979, Smalltalk\\ (Trygve Reenskaug)},4.6/MVP/{anos 1990 (Taligent);\\ comum no Android até 2017},8.1/MVVM/{2005, Microsoft (WPF);\\ padrão no Android e no Flutter},11.6/MVI/{anos 2010, inspirado\\ no Redux e no Cycle.js}}{
      \fill[ambar] (\x,0) circle (2pt);
      \node[font=\scriptsize\bfseries,text=azulpetroleo,above=1mm] at (\x,0) {\t};
      \node[rotulo,below=1mm] at (\x,0) {\d};
    }
  \end{tikzpicture}
  \raggedright
  \scriptsize
  \begin{tabularx}{\textwidth}{@{}l L L L@{}}
    \toprule
    \cabecalho{Padrão} & \cabecalho{Quem atualiza a tela} & \cabecalho{Ponto forte} & \cabecalho{Ponto fraco} \\
    \midrule
    MVC & o \emph{controller} manipula a \emph{view} & simples & no Android, a \cod{Activity} virava um \emph{controller} gigante \\
    MVP & o \emph{presenter} chama métodos de uma interface da \emph{view} & testável sem a tela & muito código de interface; guarda referência à \emph{view} \\
    MVVM & a \emph{view} observa o estado exposto pelo \emph{view model} & casa com UI declarativa e sobrevive à rotação & vários fluxos de estado podem se desencontrar \\
    MVI & estado único e imutável; cada intenção gera um novo estado & previsível e fácil de depurar & cerimônia demais para telas simples \\
    \bottomrule
  \end{tabularx}
  \normalsize
  \begin{conceito}[Na prática atual]
    Compose e Flutter usam MVVM com um estado único por tela (\cod{UiState}): uma mistura de MVVM com a disciplina do MVI.
  \end{conceito}
\end{frame}

\begin{frame}{Arquitetura limpa e arquitetura hexagonal}
  \begin{columns}[c]
    \begin{column}{0.46\textwidth}
      \centering
      \begin{tikzpicture}[every node/.style={font=\tiny,align=center}]
        \fill[cinzaclaro] (0,0) circle (2.15);
        \fill[cinzafundo,draw=azulpetroleo!40] (0,0) circle (1.6);
        \fill[ambarclaro!60,draw=ambar] (0,0) circle (1.05);
        \fill[azulpetroleo] (0,0) circle (0.55);
        \node[text=white,font=\tiny\bfseries] at (0,0) {Entidades};
        \node at (0,0.8) {Casos de uso};
        \node at (0,1.33) {Adaptadores};
        \node at (0,1.88) {Frameworks e drivers};
        \node[rotulo] at (0,-1.33) {\emph{view models}, repositórios};
        \node[rotulo] at (0,-1.88) {UI, banco, HTTP, BLE};
        \draw[setaambar] (2.6,-1.2) -- (1.2,-0.55);
        \node[rotulo,text width=1.5cm] at (3.05,-1.5) {dependências apontam para dentro};
      \end{tikzpicture}
    \end{column}
    \begin{column}{0.52\textwidth}
      \begin{definicao}[Regra da dependência]
        O código das camadas internas não conhece as externas: as regras de negócio não importam nada de Android, Flutter, HTTP ou SQL.
      \end{definicao}
      \begin{historia}
        Alistair Cockburn propôs a arquitetura hexagonal, de portas e adaptadores, em 2005. Robert C. Martin reuniu ideias parecidas na Arquitetura Limpa, em 2012.
      \end{historia}
      \begin{dica}[Na medida certa]
        Em apps pequenos, camadas de UI e de dados bastam. Crie a camada de domínio quando regras se repetirem ou precisarem de testes isolados, como recomendam o Android e o Flutter.
      \end{dica}
    \end{column}
  \end{columns}
\end{frame}

\begin{frame}[fragile]{Inversão de dependência na prática}
  \begin{columns}[T]
    \begin{column}{0.5\textwidth}
      \centering
      \begin{tikzpicture}[every node/.style={font=\tiny}]
        \node[rotulo,anchor=west] at (-1.1,0.55) {\textbf{sem inversão}};
        \node[caixaescura,font=\tiny,minimum width=2cm] (r1) at (0,0) {regras de leitura};
        \node[caixa,font=\tiny,minimum width=1.8cm] (b1) at (3.4,0) {\cod{FonteBle}};
        \draw[seta] (r1) -- node[rotulo,above]{conhece} (b1);
        \node[rotulo,anchor=west] at (-1.1,-0.75) {\textbf{com inversão}};
        \node[caixaescura,font=\tiny,minimum width=2cm] (r2) at (0,-1.3) {regras de leitura};
        \node[caixaambar,font=\tiny,minimum width=2.4cm] (i2) at (3.4,-1.3) {\emph{interface}\\ \cod{FonteLeituras}};
        \node[caixa,font=\tiny,minimum width=1.35cm] (b2) at (1.95,-2.75) {\cod{FonteBle}};
        \node[caixa,font=\tiny,minimum width=1.35cm] (m2) at (3.4,-2.75) {\cod{FonteMqtt}};
        \node[caixa,font=\tiny,minimum width=1.35cm] (f2) at (4.85,-2.75) {\cod{FonteFake}};
        \draw[seta] (r2) -- node[rotulo,above]{usa} (i2);
        \foreach \n in {b2,m2,f2}{\draw[setaambar] (\n.north) -- (\n.north |- i2.south);}
        \node[rotulo,below=1mm of m2] {implementam a interface};
      \end{tikzpicture}
\begin{codigo}[title=Dart]{Dart}
abstract interface class FonteLeituras {
  Stream<Leitura> leituras();
}
\end{codigo}
    \end{column}
    \begin{column}{0.48\textwidth}
      \begin{definicao}[Inversão de dependência]
        Módulos de alto nível não dependem de detalhes; ambos dependem de abstrações. Quem precisa define a interface; a camada de dados a implementa.
      \end{definicao}
\begin{codigo}[title=Kotlin]{Kotlin}
interface FonteLeituras {
  fun leituras(): Flow<Leitura>
}
\end{codigo}
      \begin{conceito}[O que se ganha]
        Trocar BLE por MQTT sem mexer nas regras, testar com uma fonte falsa e fazer a tela antes de o hardware existir.
      \end{conceito}
    \end{column}
  \end{columns}
\end{frame}

\begin{frame}{Injeção de dependência e localizador de serviços}
  \begin{columns}[T]
    \begin{column}{0.5\textwidth}
      \begin{definicao}[Inversão de controle]
        A classe não cria o que usa: alguém de fora monta os objetos e os entrega. É o princípio de Hollywood: ``não nos chame, nós chamamos você''.
      \end{definicao}
      \centering
      \begin{tikzpicture}[every node/.style={font=\tiny}]
        \node[caixaambarforte,font=\tiny\bfseries,minimum width=5.6cm] (raiz) at (0,0) {raiz de composição (\cod{Application}, \cod{main})};
        \node[caixa,font=\tiny,minimum width=1.7cm] (h) at (-1.9,-1) {cliente HTTP};
        \node[caixa,font=\tiny,minimum width=1.7cm] (b) at (0,-1) {banco};
        \node[caixapetroleo,font=\tiny,minimum width=1.9cm] (r) at (-0.95,-2) {repositório};
        \node[caixaescura,font=\tiny,minimum width=1.7cm] (v) at (1.9,-2) {\emph{view model}};
        \draw[setatrac] (raiz.south -| h) -- (h); \draw[setatrac] (raiz.south -| b) -- (b);
        \draw[seta] (h) -- (r); \draw[seta] (b) -- (r);
        \draw[seta] (r) -- node[rotulo,above]{injeta} (v);
        \node[rotulo] at (1.9,-1.05) {tracejado: cria};
      \end{tikzpicture}
      \raggedright
      \begin{historia}
        Martin Fowler batizou o padrão de injeção de dependência em 2004.
      \end{historia}
    \end{column}
    \begin{column}{0.48\textwidth}
      \scriptsize
      \begin{tabularx}{\linewidth}{@{}l L@{}}
        \toprule
        \cabecalho{Abordagem} & \cabecalho{Como funciona} \\
        \midrule
        Construtor & recebe tudo pronto; dependências explícitas e fáceis de trocar em testes \\
        Contêiner & uma biblioteca monta o grafo: Hilt, Koin; Provider, Riverpod \\
        Localizador & a classe pede a um registro global (\cod{get\_it}): prático, mas esconde dependências e só falha em execução \\
        \bottomrule
      \end{tabularx}
      \normalsize
      \begin{pordentro}[Tempo de vida (escopo)]
        \textbf{Único no app}: cliente HTTP, banco. \textbf{Por tela}: \emph{view model}. \textbf{Por uso}: objeto novo a cada pedido. Escopo errado causa vazamento ou estado compartilhado por engano.
      \end{pordentro}
    \end{column}
  \end{columns}
\end{frame}

\begin{frame}{Decisões de arquitetura: trocas e registros}
  \begin{columns}[T]
    \begin{column}{0.5\textwidth}
      \begin{conceito}[Toda decisão é uma troca]
        Mais camadas trazem isolamento e testabilidade, mas também mais arquivos e indireção. Pese o tamanho do time, a vida esperada do app e o custo de errar.
      \end{conceito}
      \scriptsize
      \begin{tabularx}{\linewidth}{@{}l L@{}}
        \toprule
        \cabecalho{Princípio} & \cabecalho{Significado} \\
        \midrule
        YAGNI & não construa o que ainda não é necessário \\
        KISS & prefira a solução mais simples que funciona \\
        DRY & cada conhecimento em um só lugar \\
        Regra de três & abstraia na terceira repetição, não na primeira \\
        \bottomrule
      \end{tabularx}
    \end{column}
    \begin{column}{0.48\textwidth}
      \begin{definicao}[ADR (\emph{Architecture Decision Record})]
        Texto curto, versionado com o código, que registra uma decisão: contexto, opções, decisão e consequências. Responde no futuro: ``por que fizemos assim?''.
      \end{definicao}
      \begin{exemplo}[ADR 007 --- Estado das telas]
        \textbf{Contexto}: leituras do sensor em tempo real. \textbf{Decisão}: um \cod{UiState} imutável por tela, exposto como fluxo observável. \textbf{Consequência}: telas testáveis sem interface.
      \end{exemplo}
      \begin{historia}
        Michael Nygard propôs os ADRs em 2011.
      \end{historia}
    \end{column}
  \end{columns}
\end{frame}
